% !TEX root = ../main.tex
%% ============================================================
%%  پیشگفتار، راهنمای جزوه، نمادها و سیلابس
%% ============================================================
\chapter*{دربارهٔ این جزوه}
\addcontentsline{toc}{chapter}{دربارهٔ این جزوه}
\markboth{دربارهٔ این جزوه}{}

این جزوه برای درس «آمار و احتمال مهندسی» نوشته شده و ترتیب مطالب آن \emph{دقیقاً} مطابق سیلابس درس است: هر فصل یک جلسه است. در نوشتن آن از کتاب‌های پیشرو-نیک، راس و پاپولیس (فهرست منابع در انتهای جزوه) و جزوهٔ دست‌نویس سال‌های گذشتهٔ درس استفاده شده است. هدف این بوده که هر جلسه هم \emph{از نظر علمی کامل} باشد (تعریف دقیق، قضیه و اثبات) و هم \emph{خواندنی} (شهود، مثال‌های متنوع و نکته‌های کاربردی).

\section*{راهنمای باکس‌ها}

در سراسر جزوه، هر نوع مطلب رنگ و قالب مخصوص خودش را دارد تا مرور و پیدا کردن مطالب ساده باشد:

\begin{definition}[title={تعریف \textbar\ نمونه}]{}
تعریف‌های رسمی در باکس \textbf{آبی}. اصطلاح جدید به شکل \term{فضای نمونه}{sample space} همراه معادل انگلیسی آمده است.
\end{definition}

\begin{theorem}[title={قضیه \textbar\ نمونه}]{}
قضیه‌ها، گزاره‌ها، لم‌ها و نتیجه‌ها در باکس \textbf{بنفش}. تعریف‌ها و قضیه‌های هر جلسه شمارهٔ مشترک دارند (مثلاً تعریف ۲.۱، قضیهٔ ۲.۲، \dots).
\end{theorem}

\begin{proof}
اثبات‌ها با یک خط خاکستری در کنار متن مشخص شده‌اند و با علامت $\blacksquare$ تمام می‌شوند.
\end{proof}

\begin{example}[before upper={\leavevmode{\semifont\color{excolor}مثال \textbar\ نمونه}\quad}]{}
مثال‌ها در باکس \textbf{سبز} و پاسخ آن‌ها پس از خط‌چین آمده است.
\begin{solution}
پیشنهاد می‌شود پیش از خواندن پاسخ، خودتان مسئله را حل کنید.
\end{solution}
\end{example}

\begin{intuition}شهود و تعبیرِ غیررسمیِ مفاهیم.\end{intuition}
\begin{note}نکته‌های مهم و ارتباط مطالب با جلسه‌های دیگر.\end{note}
\begin{warning}اشتباه‌هایی که دانشجویان زیاد مرتکب می‌شوند.\end{warning}
\begin{history}{}روایت‌های کوتاه تاریخی دربارهٔ شکل‌گیری مفاهیم.\end{history}
\begin{keyformula}
\[ \text{\rl{فرمول‌های کلیدی که باید به خاطر بسپارید}} \]
\end{keyformula}

در پایان هر جلسه \textbf{خلاصهٔ جلسه}، \textbf{واژه‌نامه} و \textbf{تمرین‌های پیشنهادی} آمده است. تمرین‌های دشوارتر با \hard{} مشخص شده‌اند و برای برخی تمرین‌ها جواب نهایی (بدون راه‌حل) داده شده تا بتوانید پاسخ خود را بررسی کنید.

\section*{نمادها}

\begin{center}
\renewcommand{\arraystretch}{1.25}
\begin{tabular}{cl@{\hspace{3em}}cl}
\toprule
نماد & معنا & نماد & معنا\\
\midrule
$\Omega$ & فضای نمونه & $\P(A)$ & احتمال پیشامد $A$\\
$\omega$ & یک برآمد & $\P(A\mid B)$ & احتمال $A$ به شرط $B$\\
$\F$ & فضای پیشامدها ($\sigma$-جبر) & $|A|$ & تعداد اعضای $A$\\
$A^c$ & متمم $A$ & $2^A$ & مجموعهٔ توانی $A$\\
$A\cap B$ یا $AB$ & اشتراک & $\binom nk$ & ترکیب $k$ از $n$\\
$A\cup B$ & اجتماع & $n!$ & فاکتوریل $n$\\
$A-B$ & تفاضل & $\ind_A$ & تابع نشانگر $A$\\
$\varnothing$ & مجموعهٔ تهی / پیشامد ناممکن & $\N,\Z,\Q,\R$ & مجموعه‌های عددی\\
\bottomrule
\end{tabular}
\end{center}

\clearpage
\section*{سیلابس و برنامهٔ درس}

\begin{center}
\footnotesize
\renewcommand{\arraystretch}{1.18}
\rowcolors{2}{white}{primarylight!60}
\begin{tabularx}{\linewidth}{c l X l}
\rowcolor{primary}
\leavevmode\color{white}\semifont جلسه & \leavevmode\color{white}\semifont تاریخ & \leavevmode\color{white}\semifont موضوع & \leavevmode\color{white}\semifont تمرین / کوییز\\
\rowcolor{notebg} \textbf{۱} & شنبه ۴ مهر & \textbf{مفاهیم اولیهٔ احتمال و زمینه‌های شکل‌گیری؛ مقدمات: نظریهٔ مجموعه‌ها و ترکیبیات} & تمرین ۱ (جلسه‌های ۱ تا ۴)\\
\rowcolor{notebg} \textbf{۲} & دوشنبه ۶ مهر & \textbf{فضای احتمال و اصول سه‌گانهٔ کولموگروف در حالت گسسته} & \\
\rowcolor{notebg} \textbf{۳} & شنبه ۱۱ مهر & \textbf{احتمالات شرطی، قضیهٔ بیز و قانون احتمال کل} & \\
۴ & دوشنبه ۱۳ مهر & استقلال شرطی و قانون زنجیره‌ای؛ استقلال آماری & \\
۵ & شنبه ۱۸ مهر & متغیر تصادفی گسسته، توابع \lr{PMF} و \lr{CDF}؛ متغیرهای برنولی و دوجمله‌ای & تمرین ۲ (جلسه‌های ۵ تا ۱۰)\\
۶ & دوشنبه ۲۰ مهر & متغیرهای تصادفی پواسون و هندسی (بی‌حافظگی توزیع هندسی) & \\
۷ & شنبه ۲۵ مهر & توابع توزیع و چگالی احتمال در حالت پیوسته & \\
۸ & دوشنبه ۲۷ مهر & توزیع‌های گاوسی، نمایی و سایر توزیع‌های پیوستهٔ مهم & \\
۹ & شنبه ۲ آبان & امید ریاضی و واریانس و محاسبهٔ آن‌ها برای متغیرهای معروف & \\
۱۰ & دوشنبه ۴ آبان & ادامهٔ امید ریاضی و واریانس؛ مد و میانه & کوییز ۱\\
۱۱ & شنبه ۹ آبان & گشتاورها و توابع مشخصه & تمرین ۳ (۱۱ تا ۱۵)\\
۱۲ & دوشنبه ۱۱ آبان & تابعی از یک متغیر تصادفی: توزیع، چگالی و امید ریاضی & \\
۱۳ & شنبه ۱۶ آبان & توزیع توأم متغیرهای تصادفی: توزیع دوتایی و حاشیه‌ای & \\
۱۴ & دوشنبه ۱۸ آبان & توزیع‌های شرطی و خواص آن‌ها؛ استقلال شرطی & کوییز ۲\\
۱۵ & شنبه ۲۳ آبان & امید ریاضی و واریانس شرطی و قوانین مربوط & \\
۱۶ & دوشنبه ۲۵ آبان & توزیع توأم گاوسی و خواص آن & تمرین ۴ (۱۶ تا ۲۰)\\
۱۷ & شنبه ۳۰ آبان & کوواریانس و ضریب همبستگی و خواص آن & \\
۱۸ & دوشنبه ۲ آذر & توابع یک‌بعدی و دوبعدی از دو متغیر تصادفی & کوییز ۳؛ بازهٔ میان‌ترم\\
۱۹ & شنبه ۷ آذر & ادامهٔ توابع دوبعدی از متغیرهای تصادفی & \\
۲۰ & دوشنبه ۹ آذر & قضیهٔ حد مرکزی و کاربردهای آن & \\
۲۱ & شنبه ۱۴ آذر & نامساوی‌های احتمالاتی: مارکوف و چبیشف & تمرین ۵ (۲۲ تا ۲۴)\\
۲۲ & دوشنبه ۱۶ آذر & مفاهیم آمار، نمونه‌گیری تصادفی & \\
۲۳ & شنبه ۲۱ آذر & نظریهٔ تخمین: تخمین نقطه‌ای و بازه‌ای & \\
۲۴ & دوشنبه ۲۳ آذر & تخمین بیشینهٔ درست‌نمایی (\lr{MLE}) & کوییز ۴\\
۲۵ & شنبه ۲۸ آذر & بایاس و واریانس و محاسبهٔ آن برای تخمینگرهای \lr{MLE} & \\
۲۶ & دوشنبه ۳۰ آذر & تخمین بیشینهٔ احتمال پسین (\lr{MAP}) & تمرین ۶ (۲۵ تا ۳۰)\\
۲۷ & شنبه ۵ دی & پیشین مزدوج (\lr{Conjugate Prior}) و استنتاج بیزی & \\
۲۸ & دوشنبه ۷ دی & آزمون فرض و آزمون‌های آماری: مقدار-$p$، فرض صفر & \\
۲۹ & شنبه ۱۲ دی & آزمون یک‌طرفه و دوطرفه؛ آزمون $t$ & \\
۳۰ & دوشنبه ۱۴ دی & -- & کوییز ۵\\
\rowcolor{warnbg} -- & دوشنبه ۲۸ دی & \textbf{امتحان پایانی} (ساعت ۱۵:۰۰) & \\
\end{tabularx}
\end{center}

{\small\color{graytext}ردیف‌های رنگی: جلسه‌هایی که در این نسخه از جزوه آمده‌اند.}
